\subsection{\large{\textbf{A unified approach for algorithmic complexity on some special graph classes}}}
\label{sec:poly}
Here, we see some implications of our generalized framework of relating the parameters $\XCD(G)$ and $\SC(G)$ of $G$, respectively to the parameters $k(G^*)$ and $\alpha(G^*)$ of $G^*$. In particular, we use Theorem~\ref{thm:suff_cdclique} and Proposition~\ref{pro:omega_s} as tools for deriving both positive and negative results concerning the algorithmic complexity of the problems \CDC\ and \SCP\ on some special classes of graphs. Note that the classes of graphs for which these results are applied to evaluate the time complexity of the problems \CDC\ and \SCP, may not be limited to those studied here.

 \subsubsection{\textbf{Polynomial-time algorithms:}}
 In this section, we will see a unified approach for obtaining polynomial-time algorithms (mostly with an improvement) for some special classes of graphs, using the framework of $cd$-perfectness and the auxiliary graph $G^*$.
  
\subsubsection*{\textbf{Chordal bipartite graphs}} 
The following theorem is proved in~\cite{ShaluKiruP5Free22}.
\begin{theorem}[\cite{ShaluKiruP5Free22}] \label{thm:shalu}
    Let $G$ be a $P_6$-free chordal bipartite graph. Then $\XCD(G)=\SC(G)$, and the problem \SCP\ can be solved in $O(n^3)$ time for $P_6$-free chordal bipartite graphs.
\end{theorem}

We improve and generalize the result above in Theorem~\ref{thm:chordalbip}.
First, note the following theorem proved in~\cite{DamDomChordBipGr90}.
\begin{theorem}[\cite{DamDomChordBipGr90}] \label{thm:totdom}
    The problem {\sc Total Domination} can be solved in $O(n^2)$ time for chordal bipartite graphs.  
\end{theorem}

We now have the following observations concerning the auxiliary graph of bipartite graphs and chordal bipartite graphs.
 
\begin{observation}\label{obs:bipsquare}
Let $G=(A,B,E)$ be a bipartite graph then $G^*=G^2[A]\cup G^2[B]$. 
\end{observation}
\begin{proof}
    The observation follows from the fact that for any pair of vertices $u$ and $v$ in $G$, we have $uv\in E(G^*)$ (or $d_G(u,v)=2$) if and only if either $u,v\in A$ and $\exists$ $w\in B$ such that $u,v\in N_G(w)$  or alternatively, $u,v\in B$ and $\exists$ $w\in A$ such that $u,v\in N_G(w)$. 
\end{proof}
\begin{observation} \label{obs:chordalbip}
    Let $G=(A,B,E)$ be a chordal bipartite graph. Then, $G^*$ is chordal. 
\end{observation}
\begin{proof}
    Suppose not. Since $G^*=G^2[A]\cup G^2[B]$ (by Observation~\ref{obs:bipsquare}), we have that at least one of the subgraphs $G^2[A]$ or $G^2[B]$ is not chordal. Without loss of generality, we can assume that $G^2[A]$ is not chordal (the other case is symmetric).  Let $C_k=(a_0,a_1\ldots,a_k,a_0)$ be an induced cycle in $G^2[A]$, where $k\geq 3$. Since $G$ is bipartite, and $C_k$ is an induced cycle, this is possible only if for each $i\in \{0,1,\ldots,k\}$, there exists a vertex $b_i$ in $G$ such that $N_G(b_i)\cap V(C_k) = \{a_i,a_{i+1}\}$  (mod $k+1$). This further implies that $(a_0,b_0,a_1,b_1,a_2,\ldots,a_k,b_k,a_0)$ is an induced cycle of length at least $2k$ in $G$. Since $k\geq 3$, this contradicts the fact that $G$ is a chordal bipartite graph. Hence, the observation.
\end{proof}

Now we are ready to prove Theorem~\ref{thm:chordalbip}.\\ %We then have the following theorem.} % \dhanya{due to Theorem~\ref{thm:totdom}}

\noindent\textit{Proof of Theorem~\ref{thm:chordalbip}}
    Clearly, $G$ is $\mathcal{H}$-free. By Observation~\ref{obs:chordalbip}, we have that $G^*$ is chordal, and therefore, perfect. Since any induced subgraph $H$ of $G$ is also chordal bipartite, it follows from Corollary~\ref{corr:suffcdperfect} that $G$ is  $cd$-perfect. Further, as $G$ is triangle-free, by Proposition~\ref{pro:triangle-free}, we have that $\XCD(G)=$\raisebox{2pt}{$\gamma$}$_t(G)$. Since 
 $\XCD(G)=\SC(G)$, the latter statement of the theorem is now immediate from Theorem~\ref{thm:totdom}. \qed


\begin{remark}
    Let $\mathcal{C}=\{G=(A,B,E): G$ is an $\mathcal{H}$-free bipartite graph where $G^2[A]$ and $G^2[B]$ are perfect\}. As in the proof of the above theorem, by Observation~\ref{obs:bipsquare} and Corollary~\ref{corr:suffcdperfect}, we can observe that {\sc Total Domination}, \CDC, and \SCP\  are all equivalent problems on $\mathcal{C}$. Moreover, since the problem \CC\ is polynomial-time 
 solvable for perfect graphs~\cite{GrotSchrEllipsoid81}, by Theorem~\ref{thm:suff_cdclique}, we then have that all the problems {\sc Total Domination}, \CDC, and \SCP\ can be solved in polynomial time for graphs in $\mathcal{C}$.
\end{remark}

\begin{tikzpicture}[myv/.style={circle, draw, inner sep=1.5pt}]
  \node (z) at (0,0) {};

  \node[myv] (a) at (-1,0) {};
  \node[myv] (b) at (0,0) {};
  \node[myv] (c) at (1,0) {};
  \node[myv] (d) at (2,0) {};
  \node[myv] (e) at (3,0) {};
  \node[myv] (f) at (0.5,0.5) {};
  \node[myv] (g) at (1.5,0.5) {};
 

  
  \draw (a) -- (b);
  \draw (b) -- (c);
  \draw (c) -- (d);
  \draw (d) -- (e);
  \draw (b) -- (f);
  \draw (c) -- (f);
  \draw (c) -- (g);
  \draw (d) -- (g);
\end{tikzpicture}

